\section{Summary of Gaps and Relevance to MERIT}
\label{sec:relevance-to-merit}

Earlier in this chapter, we organised the literature into three core themes. Each of these themes
exposes a gap that MERIT aims to close. We will now summarise how each of these themes motivates MERIT.

\subsection*{The Extensibility Gap}
We show in Section~\ref{sec:lit-extensibility} that 
while semantic and LLM-based matchers improve the relevance of skills, they often stay fixed in terms
of their scoring dimensions or corpus-tied coefficients. We lack recruiter-configurable and JD-driven
metrics in these approaches. In MERIT, we target extensible metric extraction and suggest weights
through TF-IDF based off the current state of the technical recruitment market.
We still allow recruiters to manually adjust the weights of the skills on the job, giving them the 
final control over the scoring system, rather than assuming a fixed set of weights.

\subsection*{The Multi-Source Verification Gap}
Section~\ref{sec:lit-multi-source} argued that sticking to 
CV-only or GitHub-only approaches can often leave claims under-verified, or do not sufficiently
provide comprehensive views of a candidate. The current GitHub-centric approaches do not fuse the 
data from candidate CV, LinkedIn and GitHub, and nor do they provide traceable conflict resolution.
MERIT, on the other hand, fuses the signals received from each of these sources together per-metric in
order to allow conflicting signals to be resolved through Bayesian evidence fusion.

\subsection*{The Transparency and Interpretability Gap}
We link the algorithmic hiring transparency deficit~\cite{bogen2018,raghavan2020} to the need
for both audit trails and source-level attribution in Section~\ref{sec:lit-transparency}.
MERIT applies Shapley values at the source level~\cite{lundberg2017unified} alongside glass-box
outputs so that recruiters can see not only the score, but also why, and which source(s) contributed to that score.

\subsection*{MERIT's Approach}
\label{subsec:merit-approach}

These gaps are closed by MERIT through three core innovations. The first is extensible, JD-driven metric configuration.
The second is multi-source fusion across CV, GitHub and LinkedIn. The third is explainable scoring with source attributions and audit trails.
To the best of our knowledge, there is no prior published system that targets technical recruitment this way, combining all three of these components
in one architecture. The three components in question are:

\begin{enumerate}[label=(\roman*)]
    \item JD-driven, configurable metric extraction with suggested priorities via TF-IDF weighting.
    \item Bayesian Beta--Bernoulli fusion across CV, GitHub and LinkedIn that separates claimed evidence from demonstrated evidence.
    \item Source-level Shapley attribution with glass-box audit trails.
\end{enumerate}

The technical components of these three innovations are well established in the literature~\cite{sparckjones1972idf, gelman2013bayesian, lundberg2017unified}.
Prior work often relaxes at least one of these constraints at a time (Sections~\ref{sec:lit-extensibility},~\ref{sec:lit-multi-source}, and~\ref{sec:lit-transparency}).
MERIT's contribution lies in the combination of the three components into a single, coherent architecture
that addresses all three gaps simultaneously.

Later, we also evaluate whether we addressed these gaps in Chapter~\ref{ch:evaluation}, 
and to what extent we were able to address them. We introduce three research questions (RQ1--RQ3) to evaluate this.
The first research question (RQ1) aims to evaluate whether we can support adaptive scoring through JD-driven, 
recruiter-configurable metrics and reliable ingestion. The next research question (RQ2) tests whether 
the introduction of multi-source fusion can improve rankings over CV-only baselines when 
public work supports or contradicts the claims made on a candidate's CV. The final research question (RQ3) 
checks whether we can improve decision calibration when given a glass-box interface to the scoring system that
features source-level attributions and audit trails.
